% =============================================================================
% Flat regions on the radius axis: theory, figures and tables.
%
% Companion to notebooks/Saddle/flat_regions_v1.ipynb. Every number is printed by
% a cell of that notebook that ran, or read from the exp15 archive it names.
%
% The theory of Section 2 is cited from report/task7-slow-saddle by the labels it
% carries there. Nothing is reproved. Section 3 is new, and it is the only place
% this file derives anything.
%
% Standalone. Own preamble, no bibliography, no \cite. The macro block mirrors
% the one near the top of _Thesis_FINAL/Article3/Survey-part3-v1.tex, through
% \providecommand, so this file can later be folded into that document.
% =============================================================================
\documentclass[11pt,a4paper]{article}

\usepackage[utf8]{inputenc}
\usepackage[british]{babel}
\usepackage{amsmath, amssymb, amsthm}
\usepackage{booktabs}
\usepackage{graphicx}
\usepackage{geometry}
\usepackage{float}
\usepackage{enumitem}
\usepackage[hidelinks]{hyperref}
\geometry{margin=2.3cm}

% ---- macros mirrored from Survey-part3-v1.tex -------------------------------
\providecommand{\defeq}{\mathrel{:=}}
\providecommand{\kbar}{\overline{\kappa}}
\providecommand{\calO}{{\mathcal{O}}}
\providecommand{\NN}{{\mathbb{N}}}
\providecommand{\abs}[1]{\left \vert #1 \right \vert}
\providecommand{\norm}[1]{\left \lVert #1 \right \rVert}
\providecommand*\mykappa[1]{ \kappa_{\scriptscriptstyle \text{#1}} }
\providecommand*\myeta[1]{ \eta_{\scriptscriptstyle \text{#1}} }
\providecommand{\Rdelta}{\textsf{R-delta}}
\providecommand{\Rstep}{\textsf{R-step}}
\providecommand{\Rdfo}{\textsf{R-DFO}}
\providecommand{\Rgrad}{\textsf{R-grad}}
\providecommand{\Rrtr}{\textsf{R-RTR}}
\providecommand{\Rgrtr}{\textsf{R-gRTR}}
\providecommand{\gone}{\gamma_1}
\providecommand{\gtwo}{\gamma_2}
\providecommand{\gthree}{\gamma_3}

\newtheorem{proposition}{Proposition}
\newtheorem{lemma}[proposition]{Lemma}
\newtheorem{corollary}[proposition]{Corollary}
\theoremstyle{definition}
\newtheorem{assumption}{Assumption}
\theoremstyle{remark}
\newtheorem{remark}{Remark}

\title{Flat regions on the radius axis\\
       \large Theory, figures and tables from \texttt{flat\_regions\_v1.ipynb}}
\author{}
\date{}

\begin{document}
\maketitle

\begin{abstract}
A radius tied to a criticality measure inherits the measure's collapse near a
saddle, and the escape time grows without bound as the saddle is made slower.
This report gathers the theory the statement rests on, cited rather than
reproved, and the seven figures and three tables of
\texttt{flat\_regions\_v1.ipynb}. The model is the true Hessian throughout, so
the stall appears with perfect curvature information. Three results are new. The
growth envelope of the escape proposition is drawn on the same axes as the run
and is never crossed. The hypothesis \texttt{as: RCA} is found to hold with
equality rather than as an inequality, $\Delta_k = \overline{\mu}\norm{g_k}$ to
eight decimal places on every iteration. And the escape-time ratio between the
two criticality measures, which the earlier prediction missed by three per cent,
is recovered to within $0.3$ per cent by an integral that accounts for the branch
switch in $\tau$.
\end{abstract}

\tableofcontents

% =============================================================================
\section{What is reported here}\label{sec: scope}
% =============================================================================

The experiment \texttt{exp15\_slow\_saddle.jl} produces the scalar tables and the
collapse. The report \texttt{task7-slow-saddle} carries the proofs. This document
reports what the notebook adds to both, which is what a batch script cannot
produce: the iterate paths, the theory drawn on the same axes as the run, and two
open items settled.

Every number below is printed by a cell of
\texttt{notebooks/Saddle/flat\_regions\_v1.ipynb} that ran, or read from the
archive \texttt{exp\_2026-08-29\_21-51-19\_slow\_saddle} that the notebook names.
The notebook executes top to bottom from a clean kernel with no errors.

% =============================================================================
\section{The theory, cited}\label{sec: theory}
% =============================================================================

\subsection{The family}

For $\varepsilon > 0$,
\begin{equation}
  f_\varepsilon(x,y) = \tfrac12 x^2
      + \varepsilon\left(\tfrac14 y^4 - \tfrac12 y^2\right),
  \qquad
  \nabla f_\varepsilon = \bigl(x,\ \varepsilon y(y^2-1)\bigr),
  \label{eqn: f}
\end{equation}
with
\begin{equation}
  \nabla^2 f_\varepsilon(0,y)
    = \operatorname{diag}\bigl(1,\ \varepsilon(3y^2-1)\bigr),
  \qquad
  \nabla^2 f_\varepsilon(0,0) = \operatorname{diag}(1,-\varepsilon).
  \label{eqn: hessian}
\end{equation}
The origin is a saddle, nonsingular for every $\varepsilon > 0$. The minimisers
sit at $(0,\pm1)$, at distance $1$ from it whatever $\varepsilon$ is, so
$\varepsilon$ tunes the escape rate without moving the geometry. That separation
is what makes the family a clean instrument.

\subsection{The four statements}

We cite these by the labels they carry in \texttt{task7-slow-saddle}. We do not
restate their proofs.

\begin{description}[leftmargin=2.6em, style=nextline]
  \item[\texttt{prop: invariant}]
    The line $\{x = 0\}$ is invariant, under truncated conjugate gradient and
    under the exact Mor\'e--Sorensen solver alike. A run started on it stays on
    it, and there $\norm{g_k} = \varepsilon\abs{y_k}(1-y_k^2) \leq
    \varepsilon\abs{y_k}$ for $\abs{y} \leq 1$.
  \item[\texttt{as: RCA}]
    The hypothesis is put on the radius rule and not on a named rule. There are
    $C > 0$ and $k_0 \geq 0$ with
    $\Delta_k \leq C \max_{k_0 \leq j \leq k}\norm{g_j}$ for all $k \geq k_0$.
  \item[\texttt{prop: escape}]
    Under \texttt{as: RCA}, with $x_0 = 0$ and $0 < \abs{y_{k_0}} \leq \tfrac12$,
    the first index with $\abs{y_k} > \tfrac12$ satisfies
    \begin{equation}
      k_{\mathrm{esc}} \;>\; k_0
        + \frac{\ln\bigl(1/(2\abs{y_{k_0}})\bigr)}{\ln(1 + C\varepsilon)}.
      \label{eqn: escape}
    \end{equation}
    The proof uses no property of the model or of the subproblem solution beyond
    $\norm{s_k} \leq \Delta_k$. That is what makes it a statement about the
    radius axis, and why it holds with the exact Hessian.
  \item[\texttt{prop: rdfo}]
    $\Rgrad$ satisfies \texttt{as: RCA} with $C = \overline{\mu}$ and $k_0 = 0$,
    immediately and with no entry condition. $\Rdfo$ satisfies it for $k > k_0$
    with $C = \gthree\zeta$, but only from an index $k_0$ with
    $\Delta_{k_0} \leq \zeta\norm{g_{k_0}}$. If no such index exists the
    proposition asserts that $\Delta_k \to 0$ geometrically and the run converges
    to the saddle.
\end{description}

The corollary \texttt{prop: budget} turns \eqref{eqn: escape} around: for any
budget $K$ there is an $\varepsilon$ making $k_{\mathrm{esc}} > K$. The stall is
therefore not a corner case to be tuned away.

\subsection{The continuous limit and its constant}

The recursion $\abs{y_{k+1}} \leq \abs{y_k} + \Delta_k$ with
$\Delta_k = C\varepsilon\,\hat\omega(y_k)$ has continuous limit
$\mathrm{d}y/\mathrm{d}k = C\varepsilon\,\hat\omega(y)$, so
\begin{equation}
  k_{\mathrm{esc}} \;\sim\;
    \frac{1}{C\varepsilon}\int_{y_0}^{1/2}\frac{\mathrm{d}y}{\hat\omega(y)} .
  \label{eqn: continuous limit}
\end{equation}
For the gradient measure $\hat\omega(y) = y(1-y^2)$ and the integral is
\texttt{eqn: exact constant} of the report,
\begin{equation}
  \int_{y_0}^{1/2}\frac{\mathrm{d}y}{y(1-y^2)}
    = \ln\frac{1}{\sqrt3\,y_0} + \calO(y_0^2).
  \label{eqn: exact constant}
\end{equation}
The $\sqrt3$ comes from the factor $1-y^2$ at the escape threshold, not from the
threshold itself.

% =============================================================================
\section{New theory: the integrand for \texorpdfstring{$\tau$}{tau}}
\label{sec: tau theory}
% =============================================================================

This section is the only place this document derives anything, and it settles the
first open item.

With $\omega_k = \tau_k = \max\{\norm{g_k}, -\lambda_k^-\}$ the earlier
prediction treated $-\lambda^-$ as the constant $\varepsilon$, giving arithmetic
growth, $k_{\mathrm{esc}} \approx 1/(2\overline{\mu}\varepsilon)$ and a ratio of
escape times of $2\ln(1/(2y_0)) = 12.43$. Two facts break that.

First, $-\lambda^-$ is not constant. On the invariant line \eqref{eqn: hessian}
gives $\lambda^- = \varepsilon(3y^2-1)$, so
\begin{equation}
  -\lambda^- = \varepsilon(1 - 3y^2),
  \label{eqn: lambda minus}
\end{equation}
which decays to zero at $y = 1/\sqrt3 = 0.577350$. The notebook confirms that
value and that the escape threshold $\tfrac12$ lies below it, so
\eqref{eqn: lambda minus} is positive throughout the range but not constant on
it.

Second, $\tau$ switches branch inside the range. The two terms are equal where
$y(1-y^2) = 1-3y^2$, and the notebook locates that root by bisection at
\begin{equation}
  y^\ast = 0.460811 \;<\; \tfrac12 .
  \label{eqn: switch}
\end{equation}
Below $y^\ast$ the curvature term dominates and above it the gradient term does,
which Table~\ref{tab: branch} shows.

\begin{table}[H]
  \centering\small
  \caption{The two terms of $\tau$ on the invariant line. The branch switch at
    $y^\ast = 0.460811$ lies inside $[y_0, \tfrac12]$, so neither term is the
    measure throughout.}
  \label{tab: branch}
  \begin{tabular}{rrrl}
    \toprule
    $y$ & $y(1-y^2)$ & $1-3y^2$ & branch \\
    \midrule
    0.00100 & 0.001000 & 0.999997 & curvature \\
    0.10000 & 0.099000 & 0.970000 & curvature \\
    0.30000 & 0.273000 & 0.730000 & curvature \\
    0.40000 & 0.336000 & 0.520000 & curvature \\
    0.46081 & 0.362959 & 0.362959 & curvature \\
    0.48000 & 0.369408 & 0.308800 & gradient  \\
    \bottomrule
  \end{tabular}
\end{table}

The correct integrand is therefore the measure itself,
\begin{equation}
  \hat\omega_\tau(y) = \max\bigl\{\,y(1-y^2),\ 1-3y^2\,\bigr\},
  \label{eqn: omega tau}
\end{equation}
and the predicted ratio of escape times is the quotient of the two integrals of
\eqref{eqn: continuous limit},
\begin{equation}
  R_{\mathrm{pred}}
    = \left.\int_{y_0}^{1/2}\frac{\mathrm{d}y}{y(1-y^2)}
      \middle/
      \int_{y_0}^{1/2}\frac{\mathrm{d}y}{\hat\omega_\tau(y)}\right. .
  \label{eqn: ratio pred}
\end{equation}

\begin{table}[H]
  \centering\small
  \caption{The two integrals of \eqref{eqn: continuous limit}, by Simpson on
    $u = \ln y$. The quadrature is validated against the closed form of
    \eqref{eqn: exact constant} before being applied to
    \eqref{eqn: omega tau}, which has none.}
  \label{tab: integrals}
  \begin{tabular}{rrrrr}
    \toprule
    $y_0$ & $I_{\mathrm{grad}}$ numerical & closed form & difference
      & $I_\tau$ numerical \\
    \midrule
    $10^{-3}$ &  6.358449 &  6.358449 & $7.11\times10^{-14}$ & 0.736439 \\
    $10^{-6}$ & 13.266204 & 13.266204 & $6.57\times10^{-14}$ & 0.737438 \\
    \bottomrule
  \end{tabular}
\end{table}

Table~\ref{tab: integrals} gives
\begin{equation}
  R_{\mathrm{pred}} = 8.6341 \ \ (y_0 = 10^{-3}),
  \qquad
  R_{\mathrm{pred}} = 17.9896 \ \ (y_0 = 10^{-6}).
  \label{eqn: ratio values}
\end{equation}
Section~\ref{sec: open one} compares these with the fitted ratios.

\begin{remark}[why $I_\tau$ barely moves with $y_0$]
  $I_{\mathrm{grad}}$ diverges logarithmically as $y_0 \to 0$ and $I_\tau$ does
  not, since $\hat\omega_\tau \to 1$ there. The predicted ratio therefore grows
  like $\ln(1/y_0)$, which is why it is $8.63$ at $y_0 = 10^{-3}$ and $17.99$ at
  $10^{-6}$ while the two $I_\tau$ differ in the fourth decimal.
\end{remark}

% =============================================================================
\section{Method}\label{sec: method}
% =============================================================================

The grid is \texttt{exp15}'s and is not redefined: $\varepsilon$ over five
decades from $10^{-1}$ to $10^{-5}$, $\overline{\mu}$ and $\zeta$ over three
decades from $10^{-2}$ to $10$, $y_0 \in \{10^{-3}, 10^{-6}\}$ and
$x_0 \in \{0, 0.8\}$, escape threshold $\tfrac12$, and at most $50\,000$
iterations. Two tolerances are used and never pooled, $10^{-14}$ in the rate
regime and $10^{-5}$ in the outcome regime.

The model is \texttt{ExactHessian} and the subproblem solver truncated conjugate
gradient throughout.

Two things are worth recording about how the notebook obtains its numbers.

\paragraph{The suite is loaded once.}
\texttt{benchmark/initialisation.jl} already includes
\texttt{exp15\_slow\_saddle.jl}, so including it again would re-evaluate its
\texttt{const}s. The notebook loads the suite once and asserts that all eighteen
names it uses came from it.

\paragraph{The archive does not carry the trajectories.}
\texttt{exp15} saves its rows through
\texttt{Base.structdiff(r, (ys = 0, deltas = 0))}, so \texttt{ys} and
\texttt{deltas} are stripped. The three tables of Section~\ref{sec: tables} read
the archive and recompute only the columns marked derived. The figures that draw
a trajectory rerun their cells: $6$ for the paths, $3$ for the envelope, $1$ for
the binding test, $14$ for the $\Rdfo$ arm and $2$ for the two measures.

% =============================================================================
\section{The figures}\label{sec: figures}
% =============================================================================

\subsection{Iterate paths}

\begin{figure}[H]\centering
  \includegraphics[width=\textwidth]{figs/flat_regions_fig1_paths.pdf}
  \caption{$\Rgrad$ capped at $\overline{\mu} = 1$ with the \emph{true} Hessian
    and truncated conjugate gradient, $y_0 = 10^{-3}$. Top row $x_0 = 0.8$,
    bottom row $x_0 = 0$, one column per $\varepsilon$. The saddle is the star
    and the two minimisers the diamonds. Trust regions are drawn on a thinned
    subset of the accepted steps. The stall is present with perfect curvature
    information.}
  \label{fig: paths}
\end{figure}

The bottom row is the invariant line, where \texttt{prop: escape} applies
exactly. The top row shows the transient: the $x$ component is solved first, and
only then does the run begin the slow climb along $y$. The panels differ only in
$\varepsilon$ and the escape iteration moves by two decades across them, while
the geometry does not move at all.

\subsection{The theory on the same axes as the run}

\begin{figure}[H]\centering
  \includegraphics[width=0.76\textwidth]{figs/flat_regions_fig2_envelope.pdf}
  \caption{$\abs{y_k}$ against $k$, both logarithmic, with the envelope
    $\abs{y_0}(1+\overline{\mu}\varepsilon)^k$ of \texttt{prop: escape} over each
    run. $\overline{\mu} = 1$, $y_0 = 10^{-3}$, true Hessian. The envelope is
    drawn up to escape only, since that is where the proposition states it.}
  \label{fig: envelope}
\end{figure}

\texttt{prop: escape} bounds the running maximum $M_k$ from above and
$\abs{y_k} \leq M_k$, so the observed curve must lie below the envelope. It does,
on every iteration of every run.

\begin{table}[H]
  \centering\small
  \caption{The envelope check. A run crossing the envelope would be a failure of
    the bound. None does.}
  \label{tab: envelope}
  \begin{tabular}{rrrrrl}
    \toprule
    $\varepsilon$ & $k_{\mathrm{esc}}$ & points
      & $\max \abs{y}/\text{envelope}$ & at $k$ & crossed \\
    \midrule
    $10^{-1}$ &   67 &   68 & 1.000000 & 0 & no \\
    $10^{-2}$ &  639 &  640 & 1.000000 & 0 & no \\
    $10^{-3}$ & 6362 & 6363 & 1.000000 & 0 & no \\
    \bottomrule
  \end{tabular}
\end{table}

The maximum ratio is attained at $k = 0$, where the two are equal by
construction. Away from the start the observed curve falls below the envelope,
which is the direction the proposition predicts.

\subsection{The radius is the binding quantity}

\begin{figure}[H]\centering
  \includegraphics[width=0.76\textwidth]{figs/flat_regions_fig3_binding.pdf}
  \caption{$\Delta_k$, $C\norm{g_k}$ with $C = \overline{\mu} = 1$, and
    $\norm{s_k}$ along one run at $\varepsilon = 10^{-2}$, $y_0 = 10^{-3}$, true
    Hessian. The three curves coincide.}
  \label{fig: binding}
\end{figure}

\texttt{prop: escape} uses only $\norm{s_k} \leq \Delta_k$ together with
\texttt{as: RCA}, so the bound is about the radius and not about the subproblem
solution. Over the $640$ iterations to escape,
\begin{equation}
  \min_k \frac{\Delta_k}{\overline{\mu}\norm{g_k}}
    = \max_k \frac{\Delta_k}{\overline{\mu}\norm{g_k}}
    = 1.00000000,
  \label{eqn: rca equality}
\end{equation}
the trust-region constraint is active on $100\%$ of iterations, and
$\norm{s_k} = \Delta_k$ to $10^{-8}$ on $100\%$.

This is stronger than \texttt{as: RCA} requires. The assumption is an inequality
and the run satisfies it as an equality, so the escape bound is not being carried
by slack in the hypothesis. The radius rule alone sets the step length.

\subsection{The \texorpdfstring{$\Rdfo$}{R-DFO} entry condition}

\begin{figure}[H]\centering
  \includegraphics[width=\textwidth]{figs/flat_regions_fig4_dfo_entry.pdf}
  \caption{$\Rdfo$ at $\varepsilon = 10^{-2}$, $y_0 = 10^{-3}$, $\zeta = 1$, true
    Hessian. Left, started inside the band at $\Delta_0 = \zeta\norm{g_0}$, so
    the entry condition of \texttt{prop: rdfo} holds at $k_0 = 0$ and the run
    stalls for $610$ iterations. Right, started at $\Delta_0 = 1$, outside the
    band, where the run never enters it. The entry index is computed from the
    trace as the first $k$ with $\Delta_k \leq \zeta\norm{g_k}$.}
  \label{fig: dfo}
\end{figure}

\begin{table}[H]
  \centering\small
  \caption{Which branch of \texttt{prop: rdfo} each cell falls in, at
    $\varepsilon = 10^{-2}$ and $y_0 = 10^{-3}$, with $\Delta_0$ swept. $k_0$ is
    computed from the trace and not assumed.}
  \label{tab: dfo scan}
  \begin{tabular}{rrlrlrr}
    \toprule
    $\zeta$ & $\Delta_0$ & $k_0$ & $k_{\mathrm{esc}}$ & status
      & $\abs{y_{\mathrm{end}}}$ & dist.\ saddle \\
    \midrule
    0.01 & $1.0\times10^{-7}$ & 0            & never & \texttt{max\_iter} & 0.18895 & 0.18895 \\
    0.01 & $1.0\times10^{-3}$ & 12           & never & \texttt{max\_iter} & 0.48241 & 0.48241 \\
    0.01 & $1.0\times10^{0}$  & never enters & 1     & \texttt{stalled}   & 1.00000 & 1.00000 \\
    \midrule
    0.1  & $1.0\times10^{-6}$ & 0            & 6048  & \texttt{stalled}   & 1.00000 & 1.00000 \\
    0.1  & $1.0\times10^{-3}$ & 9            & 5060  & \texttt{stalled}   & 1.00000 & 1.00000 \\
    0.1  & $1.0\times10^{0}$  & never enters & 1     & \texttt{stalled}   & 1.00000 & 1.00000 \\
    \midrule
    1    & $1.0\times10^{-5}$ & 0            & 610   & \texttt{stalled}   & 1.00000 & 1.00000 \\
    1    & $1.0\times10^{-3}$ & 6            & 524   & \texttt{stalled}   & 1.00000 & 1.00000 \\
    1    & $1.0\times10^{0}$  & never enters & 1     & \texttt{stalled}   & 1.00000 & 1.00000 \\
    \midrule
    10   & $1.0\times10^{-4}$ & 0            & 68    & \texttt{stalled}   & 1.00000 & 1.00000 \\
    10   & $1.0\times10^{-3}$ & 3            & 60    & \texttt{stalled}   & 1.00000 & 1.00000 \\
    10   & $1.0\times10^{0}$  & 22           & 1     & \texttt{stalled}   & 1.00000 & 1.00000 \\
    \bottomrule
  \end{tabular}
\end{table}

Table~\ref{tab: dfo scan} settles the first branch and leaves the second open.
Nine of twelve cells enter the band, at $k_0$ between $0$ and $22$, and every one
of them stalls in proportion to $\zeta\varepsilon$, which is the first branch
working. The entry index is not a formality: at $\zeta = 0.01$ the two cells that
enter differ in $k_0$ by twelve iterations and both fail to escape within
$50\,000$.

\textbf{The second branch is not occupied on this scan.} Three cells never enter
the band, and none of the three converges to the saddle. All three escape at
$k = 1$ instead, because the first step from $\Delta_0 = 1$ is enormous relative
to $\norm{g_0} \approx \varepsilon y_0 = 10^{-5}$ and carries the iterate past
the threshold in one move. The proposition's second branch requires
$\Delta_k \to 0$ geometrically, which needs a sequence of unsuccessful
iterations that does not occur here.

We report the absence rather than searching for a configuration that produces the
predicted picture. Whether the second branch is occupied anywhere on this family
is open.

\subsection{The two measures}

\begin{figure}[H]\centering
  \includegraphics[width=0.76\textwidth]{figs/flat_regions_fig5_measures.pdf}
  \caption{$\abs{y_k}$ against $k$ for $\omega_k = \norm{g_k}$ and
    $\omega_k = \tau_k$ at $\overline{\mu} = 1$, $\varepsilon = 10^{-2}$,
    $y_0 = 10^{-3}$, true Hessian, with the continuous-limit curves of
    \eqref{eqn: continuous limit} over each. The predicted curves use
    $\hat\omega(y) = y(1-y^2)$ and \eqref{eqn: omega tau} respectively.}
  \label{fig: measures}
\end{figure}

The gradient measure gives $\Delta_k \propto \abs{y_k}$ and geometric growth. The
second measure gives a radius that does not vanish with $\abs{y_k}$, so growth is
faster, and the run escapes at $k = 74$ against $k = 639$. Drawing
\eqref{eqn: continuous limit} over the observed curves is what turns the visual
comparison into a test, and Section~\ref{sec: open one} quantifies it.

\subsection{The outcome map}

\begin{figure}[H]\centering
  \includegraphics[width=\textwidth]{figs/flat_regions_fig6_outcomes.pdf}
  \caption{Three outcomes over the $(\overline{\mu},\varepsilon)$ grid on the
    invariant line, $\Rgrad$ capped with the true Hessian. Crosses mark false
    convergence, circles fast escape, squares slow escape or none. The two
    tolerance regimes are separate maps. \textbf{The first-order tolerance is
    what draws the false-convergence boundary}, $10^{-14}$ on the top row and
    $10^{-5}$ on the bottom.}
  \label{fig: outcomes}
\end{figure}

\begin{table}[H]
  \centering\small
  \caption{Outcome composition, $\Rgrad$ capped on the invariant line, twenty
    cells per row. The two regimes differ in the first-order tolerance and in
    nothing else.}
  \label{tab: outcomes}
  \begin{tabular}{llrrrr}
    \toprule
    regime & $y_0$ & false convergence & fast & slow & never \\
    \midrule
    rate, $\texttt{tol} = 10^{-14}$    & $10^{-3}$ &  0 & 0 & 10 & 10 \\
    rate, $\texttt{tol} = 10^{-14}$    & $10^{-6}$ &  0 & 0 & 10 & 10 \\
    \midrule
    outcome, $\texttt{tol} = 10^{-5}$  & $10^{-3}$ & 16 & 0 &  4 &  0 \\
    outcome, $\texttt{tol} = 10^{-5}$  & $10^{-6}$ & 20 & 0 &  0 &  0 \\
    \bottomrule
  \end{tabular}
\end{table}

The contrast is complete. At $\texttt{tol} = 10^{-14}$ not one of the forty cells
reports false convergence. At $\texttt{tol} = 10^{-5}$, thirty-six of forty do,
and at $y_0 = 10^{-6}$ every cell does. Near the saddle
$\norm{g_k} \approx \varepsilon\abs{y_k}$, so lowering $\varepsilon\abs{y_0}$
below the tolerance defeats the stopping test, and the two regimes are two
different experiments rather than two views of one.

\subsection{The gap}

\begin{figure}[H]\centering
  \includegraphics[width=0.72\textwidth]{figs/flat_regions_fig7_gap.pdf}
  \caption{The ratio of the fitted geometric rate to
    $\ln(1+\overline{\mu}\varepsilon)$, against $\overline{\mu}\varepsilon$, over
    the whole grid on the invariant line. The dashed line at $1$ is the bound.
    Every point lies on or below it.}
  \label{fig: gap}
\end{figure}

% =============================================================================
\section{The tables}\label{sec: tables}
% =============================================================================

\subsection{The \texorpdfstring{$(\overline{\mu},\varepsilon)$}{(mu-bar, eps)} grid}

\begin{table}[H]
  \centering\footnotesize
  \caption{$\Rgrad$ capped, rate regime, $x_0 = 0$, true Hessian. Read from the
    archive; $\ln(1+\overline{\mu}\varepsilon)$, the ratio and the collapse
    product are derived here. \texttt{never} is a category and never a number.
    $n$ is the point count of the rate fit.}
  \label{tab: grid grad}
  \begin{tabular}{rrrrrrrrrl}
    \toprule
    $y_0$ & $\overline{\mu}$ & $\varepsilon$ & $k_{\mathrm{esc}}$
      & fitted rate & $\ln(1+\overline{\mu}\varepsilon)$ & ratio
      & $\overline{\mu}\varepsilon k_{\mathrm{esc}}$ & $n$ & outcome \\
    \midrule
    $10^{-3}$ & 0.01 & $10^{-1}$ & 6362  & 0.000990 & 0.001000 & 0.9904 & 6.3620 & 6362  & slow \\
    $10^{-3}$ & 0.01 & $10^{-2}$ & never & 0.000100 & 0.000100 & 0.9989 &        & 50001 & never \\
    $10^{-3}$ & 0.01 & $10^{-3}$ & never & 0.000010 & 0.000010 & 1.0000 &        & 50001 & never \\
    $10^{-3}$ & 0.01 & $10^{-4}$ & never & 0.000001 & 0.000001 & 1.0000 &        & 50001 & never \\
    $10^{-3}$ & 0.01 & $10^{-5}$ & never & 0.000000 & 0.000000 & 1.0000 &        & 50001 & never \\
    $10^{-3}$ & 0.1  & $10^{-1}$ & 639   & 0.009857 & 0.009950 & 0.9906 & 6.3900 & 639   & slow \\
    $10^{-3}$ & 0.1  & $10^{-2}$ & 6362  & 0.000990 & 0.001000 & 0.9904 & 6.3620 & 6362  & slow \\
    $10^{-3}$ & 0.1  & $10^{-3}$ & never & 0.000100 & 0.000100 & 0.9989 &        & 50001 & never \\
    $10^{-3}$ & 1    & $10^{-1}$ & 67    & 0.094540 & 0.095310 & 0.9919 & 6.7000 & 67    & slow \\
    $10^{-3}$ & 1    & $10^{-2}$ & 639   & 0.009857 & 0.009950 & 0.9906 & 6.3900 & 639   & slow \\
    $10^{-3}$ & 1    & $10^{-3}$ & 6362  & 0.000990 & 0.001000 & 0.9904 & 6.3620 & 6362  & slow \\
    $10^{-3}$ & 10   & $10^{-1}$ & 12    & 0.587519 & 0.693147 & 0.8476 & 12.0000 & 12   & slow \\
    $10^{-3}$ & 10   & $10^{-2}$ & 69    & 0.094059 & 0.095310 & 0.9869 & 6.9000 & 69    & slow \\
    $10^{-3}$ & 10   & $10^{-3}$ & 642   & 0.009856 & 0.009950 & 0.9905 & 6.4200 & 642   & slow \\
    $10^{-3}$ & 10   & $10^{-4}$ & 6364  & 0.000990 & 0.001000 & 0.9904 & 6.3640 & 6364  & slow \\
    \midrule
    $10^{-6}$ & 0.01 & $10^{-1}$ & 13273 & 0.000997 & 0.001000 & 0.9976 & 13.2730 & 13273 & slow \\
    $10^{-6}$ & 0.1  & $10^{-1}$ & 1334  & 0.009926 & 0.009950 & 0.9976 & 13.3400 & 1334  & slow \\
    $10^{-6}$ & 0.1  & $10^{-2}$ & 13273 & 0.000997 & 0.001000 & 0.9976 & 13.2730 & 13273 & slow \\
    $10^{-6}$ & 1    & $10^{-1}$ & 139   & 0.095130 & 0.095310 & 0.9981 & 13.9000 & 139   & slow \\
    $10^{-6}$ & 1    & $10^{-2}$ & 1334  & 0.009926 & 0.009950 & 0.9976 & 13.3400 & 1334  & slow \\
    $10^{-6}$ & 1    & $10^{-3}$ & 13273 & 0.000997 & 0.001000 & 0.9976 & 13.2730 & 13273 & slow \\
    $10^{-6}$ & 10   & $10^{-1}$ & 22    & 0.658274 & 0.693147 & 0.9497 & 22.0000 & 22    & slow \\
    $10^{-6}$ & 10   & $10^{-2}$ & 142   & 0.094980 & 0.095310 & 0.9965 & 14.2000 & 142   & slow \\
    $10^{-6}$ & 10   & $10^{-3}$ & 1336  & 0.009927 & 0.009950 & 0.9976 & 13.3600 & 1336  & slow \\
    $10^{-6}$ & 10   & $10^{-4}$ & 13276 & 0.000997 & 0.001000 & 0.9976 & 13.2760 & 13276 & slow \\
    \bottomrule
  \end{tabular}
\end{table}

Table~\ref{tab: grid grad} shows the rows with a finite rate; the notebook prints
all forty, including three at $y_0 = 10^{-6}$ where the escape phase is empty and
the fit returns no value. The escape time is a function of the product
$\overline{\mu}\varepsilon$ alone, and halving the cap does what halving
$\varepsilon$ does. Every cell with $\overline{\mu}\varepsilon = 10^{-3}$ has
$k_{\mathrm{esc}}$ within two iterations of $6362$, whatever $\overline{\mu}$ and
$\varepsilon$ separately are.

\subsection{The \texorpdfstring{$\zeta$}{zeta} grid}

\begin{table}[H]
  \centering\footnotesize
  \caption{$\Rdfo$ started inside the band at $\Delta_0 = \zeta\norm{g_0}$, so
    the entry condition holds at $k_0 = 0$ by construction. Rate regime,
    $x_0 = 0$, true Hessian. The bound uses $C = \gthree\zeta$ with
    $\gthree = 2.0$. Rows with $k_{\mathrm{esc}} = $ \texttt{never} are rows, not
    omissions.}
  \label{tab: grid dfo}
  \begin{tabular}{rrrrrrrrl}
    \toprule
    $y_0$ & $\zeta$ & $\varepsilon$ & $k_{\mathrm{esc}}$ & fitted rate
      & $\ln(1+\gthree\zeta\varepsilon)$ & ratio
      & $\gthree\zeta\varepsilon k_{\mathrm{esc}}$ & outcome \\
    \midrule
    $10^{-3}$ & 0.01 & $10^{-1}$ & 6048  & 0.001030 & 0.001998 & 0.5154 & 12.0960 & slow \\
    $10^{-3}$ & 0.01 & $10^{-2}$ & never & 0.000104 & 0.000200 & 0.5189 &         & never \\
    $10^{-3}$ & 0.1  & $10^{-1}$ & 610   & 0.010203 & 0.019803 & 0.5152 & 12.2000 & slow \\
    $10^{-3}$ & 0.1  & $10^{-2}$ & 6048  & 0.001030 & 0.001998 & 0.5154 & 12.0960 & slow \\
    $10^{-3}$ & 1    & $10^{-1}$ & 68    & 0.093646 & 0.182322 & 0.5136 & 13.6000 & slow \\
    $10^{-3}$ & 1    & $10^{-2}$ & 610   & 0.010203 & 0.019803 & 0.5152 & 12.2000 & slow \\
    $10^{-3}$ & 1    & $10^{-3}$ & 6048  & 0.001030 & 0.001998 & 0.5154 & 12.0960 & slow \\
    $10^{-3}$ & 10   & $10^{-1}$ & 11    & 0.525981 & 1.098612 & 0.4788 & 22.0000 & slow \\
    $10^{-3}$ & 10   & $10^{-2}$ & 68    & 0.093646 & 0.182322 & 0.5136 & 13.6000 & slow \\
    $10^{-3}$ & 10   & $10^{-3}$ & 610   & 0.010203 & 0.019803 & 0.5152 & 12.2000 & slow \\
    $10^{-3}$ & 10   & $10^{-4}$ & 6048  & 0.001030 & 0.001998 & 0.5154 & 12.0960 & slow \\
    \midrule
    $10^{-6}$ & 0.01 & $10^{-1}$ & 12694 & 0.001036 & 0.001998 & 0.5188 & 25.3880 & slow \\
    $10^{-6}$ & 0.1  & $10^{-1}$ & 1280  & 0.010288 & 0.019803 & 0.5195 & 25.6000 & slow \\
    $10^{-6}$ & 0.1  & $10^{-2}$ & 12694 & 0.001036 & 0.001998 & 0.5188 & 25.3880 & slow \\
    $10^{-6}$ & 1    & $10^{-1}$ & 139   & 0.095948 & 0.182322 & 0.5263 & 27.8000 & slow \\
    $10^{-6}$ & 1    & $10^{-2}$ & 1280  & 0.010288 & 0.019803 & 0.5195 & 25.6000 & slow \\
    $10^{-6}$ & 1    & $10^{-3}$ & 12694 & 0.001036 & 0.001998 & 0.5188 & 25.3880 & slow \\
    $10^{-6}$ & 10   & $10^{-1}$ & 23    & 0.592288 & 1.098612 & 0.5391 & 46.0000 & slow \\
    $10^{-6}$ & 10   & $10^{-2}$ & 139   & 0.095948 & 0.182322 & 0.5263 & 27.8000 & slow \\
    $10^{-6}$ & 10   & $10^{-3}$ & 1280  & 0.010288 & 0.019803 & 0.5195 & 25.6000 & slow \\
    $10^{-6}$ & 10   & $10^{-4}$ & 12694 & 0.001036 & 0.001998 & 0.5188 & 25.3880 & slow \\
    \bottomrule
  \end{tabular}
\end{table}

The ratio column sits near $0.52$ across the whole grid, against near $0.99$ for
$\Rgrad$ in Table~\ref{tab: grid grad}. The $\Rdfo$ bound holds and is loose by a
factor of about two, uniformly. The reason is that $\gthree$ is the worst case of
an expansion the next iteration undoes, so the effective constant is $\zeta$
rather than $\gthree\zeta$, and $\gthree = 2$. Scaling by $\zeta$ instead puts
the collapse product at about $6.05$, close to the $\Rgrad$ value of $6.36$.

\subsection{The two measures}

\begin{table}[H]
  \centering\footnotesize
  \caption{Fitted ratio of escape times against the prediction
    \eqref{eqn: ratio values}, one row per cell that escaped under both measures.
    The fitted ratio is a quotient of two integers, so $1/k(\tau)$ is the finest
    difference it can express. That resolution is a column and not a filter.}
  \label{tab: measures}
  \begin{tabular}{rrrrrrrrr}
    \toprule
    $y_0$ & $\overline{\mu}$ & $\varepsilon$ & $k(\norm{g})$ & $k(\tau)$
      & fitted & predicted & difference & resolution \\
    \midrule
    $10^{-3}$ & 0.01 & $10^{-1}$ & 6362 & 736 &  8.6440 &  8.6341 &   0.0100 & 0.14\% \\
    $10^{-3}$ & 0.1  & $10^{-1}$ &  639 &  74 &  8.6351 &  8.6341 &   0.0011 & 1.35\% \\
    $10^{-3}$ & 0.1  & $10^{-2}$ & 6362 & 736 &  8.6440 &  8.6341 &   0.0100 & 0.14\% \\
    $10^{-3}$ & 1    & $10^{-1}$ &   67 &   7 &  9.5714 &  8.6341 &   0.9374 & 14.29\% \\
    $10^{-3}$ & 1    & $10^{-2}$ &  639 &  74 &  8.6351 &  8.6341 &   0.0011 & 1.35\% \\
    $10^{-3}$ & 1    & $10^{-3}$ & 6362 & 736 &  8.6440 &  8.6341 &   0.0100 & 0.14\% \\
    $10^{-3}$ & 10   & $10^{-1}$ &   12 &   3 &  4.0000 &  8.6341 &  -4.6341 & 33.33\% \\
    $10^{-3}$ & 10   & $10^{-2}$ &   69 &  10 &  6.9000 &  8.6341 &  -1.7341 & 10.00\% \\
    $10^{-3}$ & 10   & $10^{-3}$ &  642 &  76 &  8.4474 &  8.6341 &  -0.1867 & 1.32\% \\
    $10^{-3}$ & 10   & $10^{-4}$ & 6364 & 739 &  8.6116 &  8.6341 &  -0.0224 & 0.14\% \\
    \midrule
    $10^{-6}$ & 0.01 & $10^{-1}$ & 13273 & 737 & 18.0095 & 17.9896 &   0.0199 & 0.14\% \\
    $10^{-6}$ & 0.1  & $10^{-1}$ &  1334 &  74 & 18.0270 & 17.9896 &   0.0374 & 1.35\% \\
    $10^{-6}$ & 0.1  & $10^{-2}$ & 13273 & 737 & 18.0095 & 17.9896 &   0.0199 & 0.14\% \\
    $10^{-6}$ & 1    & $10^{-1}$ &   139 &   7 & 19.8571 & 17.9896 &   1.8675 & 14.29\% \\
    $10^{-6}$ & 1    & $10^{-2}$ &  1334 &  74 & 18.0270 & 17.9896 &   0.0374 & 1.35\% \\
    $10^{-6}$ & 1    & $10^{-3}$ & 13273 & 737 & 18.0095 & 17.9896 &   0.0199 & 0.14\% \\
    $10^{-6}$ & 10   & $10^{-1}$ &    22 &   3 &  7.3333 & 17.9896 & -10.6563 & 33.33\% \\
    $10^{-6}$ & 10   & $10^{-2}$ &   142 &  10 & 14.2000 & 17.9896 &  -3.7896 & 10.00\% \\
    $10^{-6}$ & 10   & $10^{-3}$ &  1336 &  76 & 17.5789 & 17.9896 &  -0.4106 & 1.32\% \\
    $10^{-6}$ & 10   & $10^{-4}$ & 13276 & 740 & 17.9405 & 17.9896 &  -0.0491 & 0.14\% \\
    \bottomrule
  \end{tabular}
\end{table}

% =============================================================================
\section{Open item one: the \texorpdfstring{$\tau$}{tau} constant}
\label{sec: open one}
% =============================================================================

The structural prediction, that the ratio does not depend on $\varepsilon$, holds
and was already established. What did not hold was the value. The earlier
prediction $2\ln(1/(2y_0)) = 12.43$ is wrong by nearly a half, and refining it to
$\hat\omega_\tau = 1-3y^2$ alone gives $8.362$ at $y_0 = 10^{-3}$, three per cent
below a measured $8.64$.

Equation \eqref{eqn: omega tau} closes the gap. Reading
Table~\ref{tab: measures} by resolution:

\begin{table}[H]
  \centering\small
  \caption{The comparison by population. The error tracks the resolution, which
    is what a quotient of small integers must do.}
  \label{tab: ratio summary}
  \begin{tabular}{lrrrr}
    \toprule
    population & $y_0$ & cells & fitted range & worst relative error \\
    \midrule
    every cell that escaped in both & $10^{-3}$ & 10 & 4.0000 to 9.5714 & 53.67\% \\
    every cell that escaped in both & $10^{-6}$ & 10 & 7.3333 to 19.8571 & 59.24\% \\
    \midrule
    $k(\tau) \geq 20$, resolution below 5\% & $10^{-3}$ & 7 & 8.4474 to 8.6440 & 2.16\% \\
    $k(\tau) \geq 20$, resolution below 5\% & $10^{-6}$ & 7 & 17.5789 to 18.0270 & 2.28\% \\
    \bottomrule
  \end{tabular}
\end{table}

On the four best-resolved cells, those at resolution $0.14\%$, the fitted ratios
are $8.6440$, $8.6440$, $8.6440$ and $8.6116$ against a predicted $8.6341$, and
$18.0095$ three times and $17.9405$ against $17.9896$. Those agree to between
$0.11\%$ and $0.27\%$.

\textbf{The discrepancy is explained.} The correction has two parts and both were
needed. Treating $-\lambda^-$ as constant was the larger error, and it is what
the $1-3y^2$ form already repaired. Ignoring the branch switch at $y^\ast$ was
the remainder, worth the difference between $8.362$ and $8.6341$, that is between
a three per cent gap and a two-tenths of one per cent gap.

\begin{remark}[what is left unaccounted for]
  The residual disagreement of $0.1$ to $0.3$ per cent on the best-resolved cells
  is not explained here. Two candidates are available and we do not separate
  them. \eqref{eqn: continuous limit} is a continuum approximation to a
  recursion, with an error of order $\overline{\mu}\varepsilon$ per step, and
  \eqref{eqn: exact constant} carries an $\calO(y_0^2)$ remainder. Both are of
  the right size and neither was isolated.
\end{remark}

% =============================================================================
\section{Open item two: the bound is one-sided}\label{sec: open two}
% =============================================================================

\texttt{prop: escape} bounds the escape time from below, so it bounds the growth
rate from above. The fitted rate must lie below
$\ln(1+\overline{\mu}\varepsilon)$, and the question is by how much.

\begin{table}[H]
  \centering\small
  \caption{The gap across the grid, rate regime, $x_0 = 0$. The ratio is the
    fitted geometric rate divided by $\ln(1+\overline{\mu}\varepsilon)$. Three of
    the forty cells give no fit and are excluded, leaving thirty-seven.}
  \label{tab: gap}
  \begin{tabular}{lrrr}
    \toprule
    population & cells & ratio range & above the bound \\
    \midrule
    every cell with a fit                          & 37 & 0.8476 to 1.0000 & 0 \\
    $\overline{\mu}\varepsilon \leq 10^{-1}$       & 35 & 0.9869 to 1.0000 & 0 \\
    $\overline{\mu}\varepsilon > 10^{-1}$          &  2 & 0.8476 to 0.9497 & 0 \\
    \bottomrule
  \end{tabular}
\end{table}

Two things follow.

The bound is never violated. Not one of the thirty-seven cells has a fitted rate
above $\ln(1+\overline{\mu}\varepsilon)$, which is the side
\texttt{prop: escape} predicts and the first thing a one-sided bound must
survive.

The bound is tight where it is used. Over the thirty-five cells with
$\overline{\mu}\varepsilon \leq 10^{-1}$, the ratio ranges from $0.9869$ to
$1.0000$, so the bound is within $1.4$ per cent of the observed rate over four
decades of the product. It degrades only in the two cells at
$\overline{\mu}\varepsilon \geq 1$, to $0.8476$ and $0.9497$, and the reason is
visible in Table~\ref{tab: grid grad}: there $\ln(1+\overline{\mu}\varepsilon) =
0.693147$ against $\overline{\mu}\varepsilon = 1$, so the linearisation the
continuous limit rests on has stopped applying. That is a known corner and not a
surprise.

The tightest cells are the ones that do not escape. Rows with
$k_{\mathrm{esc}} = $ \texttt{never} have $n = 50001$ points in the fit and a
ratio of $1.0000$ to four decimals. A run that never leaves the flat region
measures the growth rate more precisely than one that escapes, because it spends
its whole budget inside the regime the bound describes.

% =============================================================================
\section{Findings for the existing subsection}\label{sec: findings}
% =============================================================================

The report \texttt{task7-slow-saddle} already carries a findings section for
\texttt{Survey-part3-v1.tex}. We do not repeat it and we do not apply it. One
finding is added, from this notebook's own integrals.

\paragraph{Line 1150, the constant of the capped escape time.}
The source reads
\begin{verbatim}
    t_{\mathrm{esc}} \;\approx\; \frac{\ln(1/\abs{y_0})}{\overline{\mu}\,\varepsilon} ,
\end{verbatim}
carrying \verb|\label{eqn: escape time capped}| at line 1151. The collapse
product is compared with three candidate constants below.

\begin{table}[H]
  \centering\small
  \caption{The mean collapse product against three candidate constants, at four
    thresholds on $\overline{\mu}\varepsilon$. The threshold is not a
    convenience: the continuous limit replaces
    $\ln(1+\overline{\mu}\varepsilon)$ by $\overline{\mu}\varepsilon$, which
    fails at the one cell where that product is $1$.}
  \label{tab: collapse}
  \begin{tabular}{rrrrrrr}
    \toprule
    $y_0$ & $\overline{\mu}\varepsilon \leq$ & cells & mean product
      & residual $\ln(1/y_0)$ & $\ln(1/(2y_0))$ & $\ln(1/(\sqrt3 y_0))$ \\
    \midrule
    $10^{-3}$ & 1         & 10 &  7.0250 & \textbf{0.1172} & 0.8104 & 0.6666 \\
    $10^{-3}$ & $10^{-1}$ &  9 &  6.4722 & 0.4355 & 0.2576 & \textbf{0.1138} \\
    $10^{-3}$ & $10^{-2}$ &  7 &  6.3786 & 0.5292 & 0.1640 & \textbf{0.0201} \\
    $10^{-3}$ & $10^{-3}$ &  4 &  6.3625 & 0.5453 & 0.1479 & \textbf{0.0041} \\
    \midrule
    $10^{-6}$ & 1         & 10 & 14.3235 & \textbf{0.5080} & 1.2011 & 1.0573 \\
    $10^{-6}$ & $10^{-1}$ &  9 & 13.4706 & 0.3450 & 0.3482 & \textbf{0.2044} \\
    $10^{-6}$ & $10^{-2}$ &  7 & 13.3050 & 0.5105 & 0.1826 & \textbf{0.0388} \\
    $10^{-6}$ & $10^{-3}$ &  4 & 13.2737 & 0.5418 & 0.1514 & \textbf{0.0075} \\
    \bottomrule
  \end{tabular}
\end{table}

The residual for $\ln(1/(\sqrt3 y_0))$ falls monotonically as the limit is
approached, to $0.0041$ and $0.0075$, while the residual for $\ln(1/\abs{y_0})$
stays near $0.54$. At the tightest threshold the first is smaller by a factor of
$133$ and $72$. Only the unrestricted row favours $\ln(1/\abs{y_0})$, and that
row contains the cell at $\overline{\mu}\varepsilon = 1$ where the limit does not
apply.

This is a different location from the finding the report makes at line 1107,
which concerns the caption of \texttt{tab: slow saddle} and the model axis under
$\Rdelta$. Line 1150 sits in the radius-axis paragraph and is the equation the
collapse test at line 1157 is stated against, so it is what this notebook
measures. The same constant stands at line 1081 in \texttt{eqn: escape time},
where the source says \texttt{linearises to}, so it is offered there as an
approximation and the case is weaker.

We describe and do not apply. Both readings stand. The source may keep
$\ln(1/\abs{y_0})$ as the leading behaviour of the linearised flow, or take
$\ln(1/(\sqrt3\abs{y_0}))$ as the constant the collapse lands on. The author
chooses.

\paragraph{Nothing else is new.} The envelope, the binding radius and the two
tolerance regimes confirm statements the report already makes. The corrected
$\tau$ constant of Section~\ref{sec: open one} and the unoccupied second branch
of \texttt{prop: rdfo} in Section~\ref{sec: figures} are findings about the
report rather than about the subsection.

% =============================================================================
\section{Discussion}\label{sec: discussion}
% =============================================================================

\paragraph{The stall is a property of the radius, not of the model.}
Every run in this report uses the exact Hessian. The model sees the escape
direction perfectly and reports $\lambda_{\min} = -\varepsilon$ at the saddle.
The method still takes $6362$ iterations to move a distance $\tfrac12$ at
$\overline{\mu}\varepsilon = 10^{-3}$. \eqref{eqn: rca equality} says why: the
radius is exactly $\overline{\mu}\norm{g_k}$ on every iteration, and near the
saddle $\norm{g_k} \approx \varepsilon\abs{y_k}$, so the step the rule permits
shrinks with the distance the run is trying to cover. The subproblem solver is
not the constraint and neither is the model.

\paragraph{Two uses of second-order information, with different effects.}
The radius measure and the stopping test are both second-order and they act on
different things. Replacing $\norm{g_k}$ by $\tau_k$ in the radius shortens the
escape by a factor of $8.6$ at $y_0 = 10^{-3}$ and $18.0$ at $10^{-6}$, and
Section~\ref{sec: tau theory} now predicts both to under half a per cent. The
stopping test does something else: Table~\ref{tab: outcomes} shows that at
$\texttt{tol} = 10^{-5}$ thirty-six of forty cells stop at the saddle and report
success, and lowering the tolerance removes all of it. Neither substitutes for
the other and they are reported separately throughout.

\paragraph{The entry condition is where $\Rgrad$ and $\Rdfo$ part.}
$\Rgrad$ satisfies \texttt{as: RCA} at $k_0 = 0$ with no condition, so the stall
is unconditional. $\Rdfo$ needs the run to enter the band first, and
Table~\ref{tab: dfo scan} shows $k_0$ ranging from $0$ to $22$ and three cells
never entering. Where entry happens the stall follows, with a bound loose by a
uniform factor of about two because $\gthree$ is a worst case. Where it does not,
the behaviour on this family is immediate escape rather than the convergence to
the saddle the proposition's second branch describes. That branch remains
unexercised.

\paragraph{What the section can now claim.}
Three statements are supported by measurement rather than by argument. The growth
envelope of \texttt{prop: escape} is never crossed, on any run, at any
$\varepsilon$. The hypothesis \texttt{as: RCA} holds with equality on this
family, so the bound is not resting on slack. And the bound is tight to $1.4$ per
cent over four decades of $\overline{\mu}\varepsilon$, degrading only where
$\ln(1+t) \approx t$ stops holding. A bound that is one-sided, never violated and
tight to a few per cent is worth stating as such.

\paragraph{What remains open.}
The second branch of \texttt{prop: rdfo} is not occupied on this scan and we do
not know whether it is occupied anywhere on this family. The residual $0.1$ to
$0.3$ per cent in Section~\ref{sec: open one} has two candidate explanations and
neither was isolated. Both are recorded rather than resolved.

\end{document}
